\section{Cohorts, transitions, and recovery}\label{cohort-recovery}\label{cohorts}

A cohort is a frozen roster of model versions scored on one fresh hidden batch.
It lasts at least \(K\) subnet tempos, frozen in its manifest. The launch value
is \(K=20\). A tempo is the subnet's epoch length, 360 blocks; at about 12 seconds
per block it lasts about 72 minutes~\cite{tempo}, so a cohort takes about 24 hours
and the network completes one cohort per day. On mainnet, batch model APIs
could generate a 4,000-case batch more cheaply, but they queue requests with no
stated completion time. Until an authorized phase for generating ahead exists
(Section~\ref{protocol-checks-before-launch}), a cohort that uses them may run
longer than a day; generating at standard prices instead costs more. Rehearsal manifests may use
\(K=4\) so that deletion and recovery tests do not take days. The tempo itself
stays at 360 blocks: a longer tempo would delay every native commit--reveal by
at least a full epoch and stretch the owner's hyperparameter rate limit.

Each phase advances at a tempo boundary once its work is durable. Admission and
the transition also have frozen minimums, and the whole cohort cannot end
before its \(K\) tempos. Generation, scoring, and publication have no minimum:
they advance on completion, in order, and may share a boundary. If work takes
longer, the same cohort advances at a later boundary~\cite{tempo}; nothing is
skipped or discarded. Evaluation and publication are coordinator-owned work;
each validator's submission and finalization record is separate. Other
validators' acknowledgments, stake adoption, and software versions never gate
cohort progress. This includes the coordinator's own chain delivery: its
serverless evaluation and scheduling can continue while a weight submission
remains unresolved.

\begin{longtable}{p{0.17\linewidth}p{0.12\linewidth}p{0.33\linewidth}p{0.28\linewidth}}
\caption{Phases of a cohort with \(K=20\). Tempo positions are nominal
targets counted from the cohort's opening boundary.}\label{tab:cohort-phases}\\
\toprule
Phase & Tempos & Permitted work & Advancement \\
\midrule
\endfirsthead
\toprule
Phase & Tempos & Permitted work & Advancement \\
\midrule
\endhead
Admission & 0--1 & Accept finalized digest commitments and presigned R2
uploads of artifact files; freeze the roster. & At least one tempo. \\
Generation & 1--3 & Draw human texts under the frozen pool manifest, generate
fresh AI-side text in parallel lanes, and seal the batch and roster digests. &
On completion. \\
Scoring & 3--15 & Publish every admitted artifact as scoring starts; score
every entrant, \(R_k\), and \(R_{k-1}\) in pairs in parallel lanes, then each
executed entrant's probe run; work is sized to at most half of this window. &
On completion. \\
Publication & 15--16 & Bundle receipts; seal and sign the private score
package, the receipt index, and the public quantized result. & On completion. \\
Transition & 16--20 & Issue no new challenges. Validators process their chain
obligations independently. At the closing boundary, publish the openings and
the quantized result's values. & At least two tempos and the cohort minimum. \\
Paused & any & Retain state and receive evidence from already dispatched work.
Start no new challenge and advance no cohort. & Reconcile the saved state and
chain, then authorize resumption of the same cohort. \\
\bottomrule
\end{longtable}

The signed cohort closure records the epochs at which admission closed,
publication completed, and the cohort closed. The coordinator checks them
against the manifest's frozen admission, transition, and cohort minimums before
the cohort's target activates, and the publisher checks them again before it
lists the cohort. Every manifest, rehearsals included, freezes an
admission minimum of at least one tempo and a transition minimum of at least
two, together no more than \(K\), and a reward-intent renew-by hint of at least
\(2K\) tempos, which never pauses anything. A package finished by tempo 18 keeps the daily cadence. A later one
moves the next cohort back by whole tempos. Boundaries are finalized increments
of the chain's epoch index, not block numbers modulo 360, because epochs can
slip.
A block handler need not run at an exact block. A late or missed delivery delays
progress; it does not invalidate work, and a wall-clock timer cannot itself
advance a phase. The next cohort requires a fresh coordinator start record after
the current transition completes. Future submissions may queue without
replacing an artifact in the frozen roster. Validator delivery can continue
across later cohorts without extending the transition.

\subsection{Lagged payout}\label{lagged-payout}

When cohort \(N\)'s package produces a new signed calculation target \(T_N\), a
\emph{cohort target}, it activates at the first eligible
tempo boundary after its package and retained receipts are durable: the
boundary that closes cohort \(N\), nominally the
opening of cohort \(N+1\). Cohort \(N\) is therefore paid during the next day
while cohort \(N+1\) is scored. From activation, at every tempo, the coordinator
derives \(T_N\)'s row against a fresh finalized registration snapshot and signs
a new row record whenever the paid entries change, and each validator commits
the current signed row under a per-tempo operation identity, with native
commit--reveal enabled. That is 20 commitments per validator per day, well inside the current
rate limit, and a validator that misses a few tempos stays active. Commitments
made before activation reveal under their original rules. There is no
per-validator activation gate: no validator must prove coverage of \(T_{N-1}\)
before \(T_N\) applies.

During coordinator silence no target changes; validators keep committing the
last signed row of the active target for as long as the silence lasts, since no
reward-intent authorization expires, and cohorts are delayed,
never skipped. A new miner's first payout arrives about 25--49 hours after its
commitment, so the subnet's immunity period must exceed that interval. That is
necessary but not sufficient: when pruning a non-immune neuron would leave too
few non-immune UIDs, the chain prunes the lowest-emission immune neuron
instead~\cite{immunityperiod}, and under the burn row every miner's emission is
zero.

A package that produces no cohort target carries forward the target already
active (Section~\ref{authorization-dispatch}). Until a cohort target activates,
that is the burn row, which sends the whole pool to the owner hotkey's UID. A
launch authorization, which names no manifest or package and can authorize only
the burn row, lets validators commit it from the first tempo of the first
cohort, before any manifest is sealed. Like every reward-intent authorization it
never expires, and it is accepted only while the burn row is the active target. While
it is, each package without a cohort target also authorizes the burn row, as the target already active, and
the launch authorization covers the tempos after the package that produces the
first cohort target is sealed and before that target activates. From that
activation every validator refuses the launch authorization, at each tempo's
commitment as well as when checking it, so it is never accepted again; later
packages without a cohort target carry forward the activated target, so within
a record nothing authorizes the burn row again. Only an operator's switch to
burn issues a new launch authorization, for a fresh record
(Section~\ref{operator-decisions}).

Validators commit the burn row every tempo until
some cohort first has a gated entrant whose released score reaches the margin
cutoff; the
promotion that normally follows makes \(R_0\) the predecessor for the fallback
cutoff (Sections~\ref{reference-model} and~\ref{reward-v1}). Native
validators commit the burn row too: their adapter records an explicit,
journaled policy for submitting a row whose only positive entry is the owner
hotkey's UID, used for the burn row under the launch authorization or a burn
target, and only when the chain's weight limits and the validator's own permit
allow it; otherwise that tempo defers. At the current \texttt{rtao} pin every
other row whose only positive entries are the owner's or the validator's own UID
is held instead, as for a cohort target whose credited miners have all left:
the validator defers each such tempo, commits nothing, and leaves its previous
weights on chain. Those weights can still pay a UID's new occupant, which the
no-transfer rule below forbids, so committing such a row natively is a launch
requirement (Section~\ref{protocol-checks-before-launch}).
That regime has no bounded duration; no miner earns
during it, and the burn withholds the subnet's TAO-side emission, while under
the documented rules validators and the owner still receive their
participant-side alpha (Section~\ref{unallocated-rewards}). Once the subnet's UID slots are full, a new registration can therefore replace any miner past its
immunity period; the credit of its entries still awaiting payout then burns
(Section~\ref{reward-v1}), and re-registering starts a new registration
generation.

\subsection{Fan-out evaluation work}\label{fan-out}

Bulk evaluation work runs as fan-out jobs: challenge generation, paired scoring
and probe runs, receipt bundling, and the per-model summaries that packaging
needs. The cohort's Durable Object remains the only authority. It records each
lane's assignment and holds the task leases, and lanes claim and commit work
through short, idempotent requests to it. Each lane runs in its own Durable
Object invocation with its own platform limits, and each scoring lane drives its
own Container instance. Lanes number at least one per scored artifact, and more
as the batch grows, within the account's Container limits~\cite{cfcontainerlimits},
the providers' model rate limits through one shared limiter~\cite{cfailimits},
and, on testnet, the cohort's budget. Lane count changes elapsed time and per-lane start
cost, never a document's outcome: the scorer's outputs are bitwise identical
across instance sizes (Section~\ref{existing-evidence-and-implementation-status}).
Lanes never sign; signing stays in the cohort object's phase transitions.

Containers start on demand, and nothing runs warm at rest; scoring instances may
start just in time during the tail of generation, and cold starts are accepted.
A Container crash or out-of-memory failure counts against an artifact only when
it reproduces in single-document isolation
(Section~\ref{aggregation-and-failure-policy}). The testnet roster is capped at
40 entrants. Mainnet has no roster cap beyond the subnet's 256 UIDs and no
spending cap. Providers publish budget alerts but no hard spending
cap, so on testnet lane caps and per-cohort budgets bound spending.

\subsection{Block notifications and finalized recovery}\label{block-notifications}

Live \texttt{rtao-rpc} webhooks on the finalized stream deliver every finalized
block, each to its own Worker invocation. That delivery is the coordinator's
primary heartbeat and clock. Deliveries are signed, at-least-once, and retried
within a bounded provider journal. The provider retries a delivery that receives
no success response ten times, with backoff, over about 5 hours 49 minutes, and
then pauses the subscription; nothing more is delivered until it is resumed. The
coordinator's handling of authentication, retries, ordering, and retention still
requires qualification.
Authenticate ingress and persist the requested work before acknowledging
delivery. Duplicate, delayed, or reordered notifications must not repeat a
transition, sign another attempt, reset a deadline, or skip a cohort. A
notification neither authorizes an action nor establishes that its block is
finalized.

Retain a finalized cursor with height and hash, plus completed event identities,
in Durable Object storage. Reconcile chain history through RPC, checking hash
continuity and processing missing finalized blocks in order. Advance the cursor
only after the dependent state changes are durable. Best-head notifications may
be reorganized; they cannot advance finalized state. Any conflict with the saved
finalized history pauses affected actions for diagnosis. After the cursor
advances, an invocation claims due work through a durable lease with an
idempotent operation identity, holder, and expiry; overlapping invocations
skip claimed work, and later blocks re-claim stale claims and continue
resumable steps.

A Cron Trigger~\cite{cfcron} every minute is the only fallback and the only
timer. It is a no-op while valid deliveries arrive. If none arrives for a
configurable threshold, five minutes by default, or the cursor lags the
finalized head by more than that threshold, the coordinator enters a durable
fallback mode. Each tick then reads the finalized head from public OpenTensor
RPC and processes blocks from the cursor through the same per-block path and
deduplication, bounded per tick within public rate limits. Delivery resumption
returns it to primary mode automatically. A paused subscription keeps the
coordinator in fallback mode: the watchdog reports the fallback state and the
time since the last valid delivery for an operator alert, and an operator
resumes the subscription. Automatic resumption is not yet implemented, because it needs a
subscription-management credential in the Worker or a provider resume path; until
then a pause lasts until an operator acts, a stated limit. The nominal wait is one
finalized block plus delivery latency, or up to the threshold plus one tick in
fallback.
Actual observation delay also includes scheduler delay, RPC and finality
latency, and catch-up work; this is no hard recovery bound during an outage.
Bound each invocation and persist its continuation, so a long outage needs no
resident subscriber or unbounded catch-up request.
Resume scheduling at a later eligible tempo boundary after recovery, preserving
the current cohort and its remaining active-time budgets.
Pending chain obligations have separate recovery records; advancing the block
cursor does not require their completion.

\subsection{Authorization and durable dispatch}\label{authorization-dispatch}

Coordinator-owned deadlines use a durable active-time clock. Bounded progress
intervals advance that clock only when recorded in an authoritative checkpoint;
missed intervals are not charged retroactively after recovery. Silence therefore
preserves the remaining budgets for phase timing and coordinator-set submission
windows without requiring a pause write during the outage. Recovery publishes
the remaining budgets and revised eligible tempo boundaries. Extending an open
window cannot reopen a frozen roster or a finalized outcome.

Permission to act is separate from that clock. New cohort work requires a
signed authorization bound to cohort, phase, and checkpoint generation. It has
no expiry: it stays valid until a newer generation supersedes it or a signed
revocation withdraws it, and it fences coordinator work only. Weight commitments use a reward-intent
authorization instead, a separate record signed under the authorization key.
Each package is sealed with one that names the cohort's manifest and package by
digest and the target active after it: the package's new target, whose
predecessor must be the target active when the package is sealed, or, when the
cohort produced none, the target already active. Validators receive the
authorization, published beside the coordinator's row records, but never the
package. Until the first cohort target
activates, a launch authorization covers the burn row
(Section~\ref{lagged-payout}). No reward-intent authorization expires: it yields commitments
from its issue epoch for as long as its target is the validator's active
target, and stops only when another target activates or a signed revocation
withdraws it. A coordinator outage of any length therefore leaves validators
committing the last signed row. Each record carries a renew-by hint, frozen in
the manifest at no less than \(2K\) tempos, that never refuses or pauses
anything; a renewal only refreshes that hint, and none is required. Revocation
is explicit: a signed revocation names a target and withdraws every
authorization of it issued at or before its epoch, and every validator journals
the revocations it has checked.

Validators verify each authorization, the coordinator's row record under it,
and the revocation set under the pinned authorization key before using them.
They check that the row names the authorization and a target that is their
active one or activates later, and that the authorization's issue epoch is no
earlier than the cohort's opening. Each validator, the coordinator included,
keeps its own durable record of the target it holds active, written only
forward. A validator other than the coordinator accepts a row only for that
target, with the same activation epoch, or for a target whose activation epoch,
as the signed row records it, is later than the held target's, and then records
the row's target as active. One that missed several activations therefore moves
straight to the current row's target, without the rows or packages in between.
The coordinator moves its record only through each cohort's checked closure,
so each target it activates names the target it held as its predecessor, except
at an operator's restart (Section~\ref{operator-decisions}). A launch
authorization is accepted only while no cohort target of its record has
activated: the first one before any cohort target, or a restart's, issued at the
restart floor after the target the validator holds. At every
tempo a new commitment also requires the authorized target to be the one the
validator holds active and the epoch to be at or after the authorization's
issue, so an authorization stops yielding commitments once another target
activates. Existing chain obligations retain their recorded authority and
native deadlines. If a work
authorization lapses or its checkpoint cannot be verified, no new coordinator
work begins until the live coordinator issues another; a lapse never ends or
skips a cohort and never touches validators' commitments. A fresh
finalized-chain view is required; a cached height cannot keep a work
authorization valid indefinitely. One non-coordinator validator's silence never
stops the shared clock.

Before dependent actions, retain the manifest, roster, artifact and credited
identity bindings, task and attempt identities, completed and outstanding work,
remaining active-time budgets, scores, and chain receipts. Keep referenced
artifacts and incomplete-cohort records without automatic expiry. Durable
recovery exports and restore rehearsals must cover pauses longer than any
provider backup window.

A lane's assignment is persisted with its state transition before its wakeup is
sent. The lane claims each task under a lease, checks current authorization and
generation, then retains its result before releasing the claim. Wakeups are
notifications; authoritative records reconstruct dispatch after redelivery,
loss, or retention expiry. The coordinator's state does not depend on a
surviving message, lane invocation, Worker process, or always-on machine.

Task leases and compute attempts have separate identities. Lease expiry fences
an old worker from new actions but does not erase its evidence. Reconcile an
uncertain external call before starting another attempt. Reclaiming a task or
recording an already durable answer consumes no compute retry. Starting another
computation does. A frozen attempt-order rule selects one canonical outcome;
never choose the best answer returned across retries. The late result of an
old worker cannot advance a newer checkpoint or reopen a finalized score.

\subsection{What pauses and what continues}\label{pause-behavior}

Coordinator downtime suspends protocol progress and cohort scheduling. It does
not remove unfinished tasks from the denominator, turn them into miner failures,
reroll challenges, or start a replacement cohort. Already dispatched work may
finish and persist evidence, but its completion cannot authorize another task.
Resume only after reconciling in-flight attempts and late results. If scoring
was final before the outage, resume publication or the transition without
rescoring. Missing evidence leaves the cohort paused until a reviewed resolution
is possible or the operator abandons it (Section~\ref{operator-decisions}).

A per-attempt compute limit remains the bound frozen for that attempt. Queue
delay, operator suspension, and unresolved infrastructure faults are not measured
model compute. A limit violation already established remains an outcome.
Customer delivery/refund terms, credential expiry, transaction mortality, and
native chain reveal windows retain their own clocks. Extending a
coordinator-owned deadline cannot extend those obligations or renew an expired credential.

\subsection{Operator decisions}\label{operator-decisions}

No precondition in the evaluation service, the coordinator, the publisher, or a
validator may lock the operator out. Each has a journaled operator path past
it, so the operator can always decide to start a new cohort, change the
references, install or pause a cohort template, or restart the record without
first changing code. Moving forward never requires reading retired or
superseded state, and a paused cohort calendar stays paused, across a template
install, until the operator resumes it. These paths keep the protocol's other
rules: no cohort is skipped or expired because time passed, cohort numbers are
never reused, each validator signs at most one commitment per chain, subnet
registration, hotkey, and epoch and never resubmits blindly, validators keep
committing the last signed row, and no decision recreates a lock it removes.

Each decision is journaled before it takes effect and leaves a signed,
content-free public record that states its reason. Decisions that decline a
result, change reward inputs, or restart the record are also signed with a
separate operator-decision key, kept offline with a recorded backup. That key is
distinct from the coordinator's hotkey and from its authorization and receipt
keys, and only the coordinator and the evaluation service verify it; validators
keep verifying rows, authorizations, and revocations under the pinned
authorization key alone. Rotating the operator-decision key is a recorded
configuration change of those two services that needs no signature under the
old key, so losing it locks nothing. The paths are these:
\begin{itemize}\tightlist
\item \emph{Abandon.} The operator may end any cohort that has not closed,
including a queued cohort or one with no fault. The abandonment record names the
cohort's phase and digests, the cohort produces no target, and its number is
never reused. An abandonment in packaging or the transition, when the outcome
may already be known, declines that outcome and needs the operator-decision key.
\item \emph{Refusal to pay.} The operator may refuse to pay a finished cohort by
declining its closure, by abandoning it late, or by naming it in a switch to
burn, each under the operator-decision key. A declined result is never
installed. Later targets chain from each cohort's target, so continuing after a
refusal goes through a switch to burn and a fresh record.
\item \emph{Start over.} To stop paying the current record, the operator switches
the coordinator to burn. With signing paused, the coordinator sets a restart
floor one epoch after the last finalized one, installs the burn row from the
floor under a launch authorization issued at that epoch, and names in the
signed restart record every cohort that opened before the floor without being
handled, so no cohort is dropped silently. The operator then opens a fresh
record whose first reference is \(R_0\), with no predecessor. Winners of the
earlier record are never paid again. Validators accept the restart through that
launch authorization, which was issued after the target they hold, and refuse
rows from before its floor; none needs a manual reset.
\item \emph{Reward-rule change.} A new reward rule takes effect only through a
fresh record under that rule, reached by the same switch to burn; no record
carries across reward rules.
\item \emph{Reference decision.} The operator may set the current and previous
references directly, with a recorded reason, outside a fresh record
(Section~\ref{reference-model}).
\item \emph{Re-registration.} If the subnet is registered again on the same
chain, the restart also rotates the authorization key, so no authorization or
row signed for the earlier registration verifies under the new one; validators
receive the new key through a signed release or trust update.
\end{itemize}

These paths remain available on mainnet, with the same signatures and public
records. No outage triggers one: coordinator silence pauses the cohort
(Section~\ref{pause-behavior}), and only an explicit decision ends, declines, or
restarts. A decision changes only future work and payment. Commitments already
submitted keep their recorded authority and reveal under their original rules,
and rewards already paid are never reassigned. The decisions extend the trust
the pilot already places in the operator
(Section~\ref{trust-assumptions-and-adversarial-behavior}).

\subsection{Coordinator rows and commit--reveal}\label{independent-weight-calculation}

The coordinator alone measures entrants, computes the reward outcome, and
derives the native weight row. Its private score package
(Section~\ref{aggregation-and-failure-policy}) never leaves its storage; its
public quantized result carries everything needed to recompute the row.
Once a tempo, from its own finalized registration snapshot and independently of
its own commitment, the coordinator derives the active target's row and signs a
new row record whenever the paid entries change. The record binds the
deployment and chain, the reward rule, the cohort with its manifest digest and
public result, the target and the authorization it is committed under, its
epochs and derivation block, the owner hotkey, and, per paid UID, its hotkey,
registration block, and integer weight. No chain runtime version is bound. The
record is published off chain beside the authorization and the signed
revocation set as soon as it is signed, before the coordinator's own hotkey
signs anything, and no validator waits for the coordinator's commitment. The
calculation fixes numeric precision, ordering, tie handling, and native
quantization, so the same public result and registration bindings produce
identical canonical UID/weight entries. Held-out texts, labels, per-document
predictions, exact scores, receipts, customer content, and unpublished model
weights stay coordinator-private, except that held-out passages go to the
hosted generation providers of Section~\ref{cohort-batch}. This disclosure boundary has the scoring
trust assumption stated in
Section~\ref{trust-assumptions-and-adversarial-behavior}. No zero-knowledge
proof is required.

Each other validator verifies the row record, its authorization, and the
revocation set under its pinned keys, checks that the record names the
authorization and a target it holds active or that activates later, journals
the revocations, and commits the record's entries. Its only change to the row uses its own finalized
snapshot: the units the row burns through the coordinator's recorded owner
hotkey, and the units of any registration that snapshot no longer holds, go to
the UID of that snapshot's owner hotkey, which burns them. When units must go
there and that hotkey holds no UID, the validator defers that tempo's
commitment; it never sends them to another UID. It needs nothing else to commit: it receives no score package, receipt,
document, label, or exact score, recomputes nothing in order to commit, and
downloads no model. Apart from the owner-UID deferral above, which skips only
that tempo, two events about the row pause its new commitments: a signed
revocation of its authorization, which it journals, until a newer authorization
is installed; and an installed record whose encoding its release does not
support, until a supported record is installed. A newly published record that
fails the verification above, including one that names an earlier target,
raises an alert and is never installed, so the validator keeps committing its last verified row while no
valid newer row arrives. The commit--reveal flag check below, journal
recovery, and the migrations of Section~\ref{participant-updates} can also pause
it, for reasons unrelated to the row. Under the off-chain contract, a conforming
validator does not invent a score, omit an entry, or apply an undeclared
weighting rule; native consensus does not enforce it. Frozen reward policy remains
common to all validators.

Anyone, a validator included, can recompute the row from the public quantized
result and compare it with the signed row and the revealed weights. That checks
the coordinator's reward arithmetic, not hidden scoring, and its result is
advisory: a mismatch it finds never changes the row a validator commits and
never pauses its commitments. The mismatch can be published as evidence; the
coordinator corrects a defect it confirms only through a signed revocation or a
newer signed row, which validators verify as above. A validator may also
check models in its own way and publish its own scores, equally advisory. The
operator prefers to keep native weight commit--reveal enabled, viewing its use
as a positive signal within Bittensor and accepting the reveal delay. It adds
no copy protection, since the signed row is published off chain before any
reveal.
The enabled flag check below enforces that configured submission policy.
The recognized subnet owner sets
\texttt{commit\_reveal\_weights\_enabled = true}, through the agreed
owner-authorized execution path under a lease; the operator verifies the
effective setting before live submission
\cite{weightsubmissionmode,commitreveal}. Validators check that setting against
finalized state at startup and fresh submission state before signing. An
unexpected disabled flag pauses new weight commitments rather than silently
selecting plaintext submission. Model-artifact digest commitments are separate
admission records.

Validators, including the serverless coordinator, commit the coordinator's
signed row under their own keys, nonces, journals, and timelock
material. They target each eligible native commit window from the published
schedule without waiting for another validator to commit, reveal, or finalize.
Use the qualified native timelocked path, whose reveal is performed by the chain.
A submitted commitment may therefore reveal during a coordinator outage or
before a later audit completes. The client cannot assume it can cancel that
reveal by withholding a second transaction.

Parallel commitments avoid serial reveal stages between validators; they do not
remove native reveal latency or guarantee a same-tempo or atomic switch. A missed
window defers only the affected validator's delivery. Reconcile its old attempts
and current authorized target before scheduling another eligible window. Native
submission and reveal clocks do not extend when coordinator deadlines pause.

The coordinator seals each complete evaluation and durably publishes its score
package and calculation target. Cohort completion records publication, not
validator adoption or miner payment. At the next permitted boundary after the
transition minimum, the coordinator may start another cohort even if every
external validator is offline, outdated, or refusing to submit. Evaluation,
service, updates, and recovery require no other validator's acknowledgment or
cooperation. A newer target supersedes the active one only when the coordinator
holds its complete authorized package and a verified receipt index, and only at
its activation boundary. This activation needs no coordinator chain
transaction. Revocation and identity changes still apply; the authorization
itself never lapses.

After its own row is revealed and finalized, the coordinator publishes an
observation certificate binding the score-package digest, coordinator identity
and hotkey generation, canonical row, and finalized block reference. Anyone can
compare this row and observed peer rows with the signed row records.
Coordinator finalization is a later consistency
audit, never a prerequisite for another validator's commitment. A missing or
incorrect coordinator chain row is a delivery discrepancy; it cannot invalidate
an otherwise valid signed row record or stop other validators. A shared row or
encoding defect pauses another validator's new commitments only through the
two row events listed above, a signed revocation or an installed encoding its
release does not support, while correction and recovery proceed. Retain and reconcile commitments already
submitted. Authorized coordinator hotkey rotation updates the chain identity
while retaining score-signing authority and existing obligations under their
recorded keys; it imposes no adoption gate. An owner-hotkey rotation needs no
new target either: validators read the owner hotkey from finalized state at each
tempo (Section~\ref{unallocated-rewards}).

Commit--reveal hides native submissions until reveal. The coordinator's row is
published off chain as soon as it is signed, before any reveal, and also
follows from the public quantized result, so this design does not claim to
prevent copying it or to prove that any validator checked it. Native consensus
still accepts each validator's own submission and does not enforce this
off-chain contract.

Each validator retains the signed row records and revocations it installed, its
signed attempts, commit/reveal material, and finalized outcomes. On return it reconciles
every submitted or uncertain operation before acting on a newer target. A target
never submitted may be superseded with its history retained; a new package cannot
erase an outstanding commitment or justify blindly replaying an obsolete row.
Catch-up must prevent an older pending reveal from silently overwriting a newer
row. Submission-mode migrations preserve those obligations and use reconciliation
before new signing. These constraints apply to that validator's delivery path,
never to cohort advancement or another participant.

Subtensor keeps producing blocks and applying emissions, activity, and
registration rules while this protocol is paused. Submitted validator rows do
not switch atomically. The network depends on timely validator submission for
the intended weight rows to influence native consensus. Missing or stale rows
can change actual emissions even though they cannot delay application progress.
Publish observed delivery and adoption separately from cohort completion.
SlopNinja cannot force a row, freeze native emissions, or promise that a missed
reward period will be paid later.

Bind each award to its original hotkey and registration generation, not only a
UID. If that UID changes occupant, do not transfer the old credit; the share of
a hotkey that is no longer registered under that generation burns through the
owner hotkey's UID at each tempo. Preserving the cohort grants neither suspension
of native emissions nor a right to replay or reassign rewards already paid.
Retaining all scores and delivery records provides accounting continuity, not
guaranteed deferred payment.

\subsection{R2 inference receipts and retention}\label{cohort-receipt-retention}

Each inference attempt has a content-free receipt in R2, including benchmark
and probe calls, paid-service calls, retries, failures, and uncertain outcomes.
The one exception is the coordinator's pre-flight, outage, and admission
probes: fixed synthetic calls, such as one generation or one scoring run, that
check whether a provider or the scorer answers, for example before a paused
cohort resumes by itself. They carry no evaluation content, their cost is
counted with the cohort, and they write no receipt. Receipts bind
the deployment, immutable invocation and attempt identities, model/version,
status, usage, and private-content commitments. They expose no prompt, response,
held-out label, parsed feature vector, or commitment opening. Receipts are
written in bounded chunk objects. A cohort package indexes each chunk by object
identity and digest and each receipt by its position; its immutable manifest
fixes the cutoff, count, and commitment root. The coordinator's private score
package binds that manifest. Receipts are coordinator-private audit evidence:
they are never published or sent to validators, because an attempt's outcome
class, which an artifact can steer, can reveal information about hidden
documents. The coordinator installs a new target only after it verifies that
cohort's receipt index, counts, and object digests. Receipt verification checks
the package's integrity, not honest inference or hidden scoring. Cloudflare does
not sign these assertions.

Evaluation receipts belong to their cohort. Paid-service receipt records enter
the next open certification package under a fixed cutoff rule. An uncertain
attempt remains explicitly pending; a later outcome is an appended update in a
later package under the same invocation identity, never a rewrite of sealed
evidence or another credited inference. Pending paid-service records cannot
relax the requirement to account for every assigned evaluation outcome.

Detailed receipts and their R2 index for cohort \(N\) are deleted once cohorts
\(N+1\) and \(N+2\) complete and, if cohort \(N\) produced a target \(T_N\), \(T_N\)
has been superseded; a cohort that produces no new target has no target to
supersede, so its receipts are deleted once \(N+1\) and \(N+2\) complete.
Completion means the
coordinator's durable cohort completion, independent of any validator's weight
adoption. This rolling window counts completed cohorts and superseded targets,
not wall time or elapsed tempos; a coordinator pause cannot advance deletion
eligibility, and a target kept live through cohorts that produce no new target
keeps its receipts. Published entrant artifacts stay public and are never
recalled; their admission copies are deleted once the cohort closes and
publication is recorded, and the artifacts of \(R_k\) and \(R_{k-1}\) remain
hosted for evaluation while they are references. Probe fingerprints are private evaluation records, retained
with the probe while their probe version is in use; the published-artifact index
of content digests is retained indefinitely. A durable serverless cleanup task
deletes eligible objects,
invalidates access, and records completion. Delete managed replicas and caches
as well; do not retain a separate indefinite archive of these detailed receipts.
Previously downloaded copies cannot be recalled. The testnet evaluation service
runs this deletion. One gap remains as of November 9: the receipts of the
cohorts run just before the testnet record restart
(Section~\ref{testnet-r1}) are not yet deleted, because the step that lets the
new record release them waits on an older record's store that no longer
answers.

Two completed successors are a deliberate pilot retention choice: they keep
recent audit evidence available while limiting the number of completed cohorts
whose detailed metadata remains hosted. With daily cohorts, receipts become
eligible for deletion at least about 48 hours after their cohort completes.
They guarantee neither a minimum audit duration nor a byte limit, particularly
during a coordinator pause. Roots and counts cannot reconstruct deleted
receipts, which limits later investigations.

Retain the signed score package and calculation inputs, the public result,
package root and counts, canonical row, signed row records, cohort completion
record, any coordinator observation certificate, and chain/accounting recovery
records. They remain after detailed receipts are deleted. A returning validator
uses the current signed row for new commitments and its durable journal for
earlier submitted obligations; neither needs receipts. Missing or corrupt
receipts within the retention window stop only the coordinator from installing
that cohort's target; validators keep committing the last signed row. Customer content, delivery, and refund records retain their separately
disclosed policies; this window grants no extra content retention.

\subsection{Runtime upgrades and chain intents}

Freeze evaluation meaning separately from operational builds. The models,
batch, labels, scoring rules, and numerical tolerances define the comparison.
The current RPC connection, coordinator build, and Finney runtime are recorded
execution facts; they need not remain identical for the cohort to finish. An
operational fix can preserve the frozen comparison.

No component gates on a runtime version, and no cohort, package, row, or journal
is certified for one; a cohort survives any number of upgrades unchanged.
Validators and the coordinator keep running through every upgrade. Use
finalized history for decisions and reconciliation. Before submission, refresh
runtime metadata and signing parameters, and read current submission
constraints and epoch state from one fresh best-head snapshot. Resolve each
call, storage item, signed extension, and runtime API by name in that metadata
and type-check the supplied values against the live types; repeat at signing,
since finality can lag an upgrade. A mechanically safe difference, such as a
wider integer or a new trailing optional argument, is encoded adaptively and
recorded. An operation with no type-correct encoding is deferred alone with a
durable diagnostic and re-evaluated on every new block, while ingestion,
journaling, recovery, reconciliation, scoring, and package verification
continue. Never submit a payload that does not type-check against live metadata.
A runtime version is recorded only as an observation; signing reads it fresh.
The native adapter follows this rule at its current \texttt{rtao} pin, except that
upstream still composes weight calls and reads some storage through compiled
descriptors; the ``Implementation boundary'' paragraph in
Section~\ref{native-implementation-boundary} records how those cases defer.

Persist a logical chain intent and its signed attempts before sending bytes.
Retain the cohort/result revision, participant binding, intended effect, nonce,
observed runtime, transaction hash, and commit/reveal material. A timeout means
an unknown outcome. Reconcile finalized execution, pending transactions, nonce
state, and reveal obligations before signing another attempt. New bytes required
by a changed runtime, nonce, or mortal lifetime remain attempts at the same
intent. They do not authorize a duplicate effect. Decode historical evidence
using metadata for its execution context.

The native adapter integrates \texttt{rtao}'s durable journal, scheduling,
semantic checks, and signer-free finalized recovery~\cite{rtao,nativeadapter}.
Opening it and beginning external recovery persist pause; resumption is explicit.
A successful recovery call may leave unresolved work. Inspect per-operation
outcomes and retained reveal evidence before claiming completion. Current
supported operations and implementation gaps appear in the ``Implementation
boundary'' paragraph in Section~\ref{native-implementation-boundary}.
Adapter updates must retain the original cohort and journal.
The serverless coordinator needs equivalent durable intent and attempt records
in Durable Objects, exclusive signing ownership, and qualified Workers transport
and cryptography. A native file journal or passing native test does not establish
this integration. Recover the coordinator's submitted obligations before new
signing without making that recovery a barrier to unrelated cohort work.

\subsection{Automatic updates and resumable migrations}\label{participant-updates}

Automatic updating is selected for the pilot. Whether participants may opt out
or defer installation, and how releases are staggered across installations,
remain deployment-policy decisions to settle before launch. If local deferral
is permitted, it cannot bypass signing compatibility or pending-obligation
recovery; rollout sequencing cannot require other validators' acknowledgments
or gate cohort progress on their adoption.

One Rust supervisor per native installation stages signed releases and owns
process handoff and restart. Miner and validator applications automatically
fetch and verify release manifests and matching platform artifacts. The manifest
binds artifact digests, protocol compatibility, schema versions, migration
order, and release-key generation. The updater rejects unauthorized or
superseded release metadata under the published key-rotation and rollback
policy. Each participant updates independently; neither installing a release
nor resuming afterward requires other validators to acknowledge it. An exclusive
journal lock and signer ownership permit only one active participant process to
write that journal or sign with its hotkey; the supervisor fences the old
process before handoff.

Retain ordered migration paths from every released persistent schema. A
participant that missed many releases applies every required migration in
order, even if intermediate application binaries need not run. Each step has a
durable progress record, validates its input and output, and resumes safely
after interruption.
Pause local new work and signing, fence the old process, and preserve a recoverable
copy before changing state. Never replace a journal with an empty one, discard
pending commitments, or blindly restart an old binary against a migrated schema.
After migration, recover finalized chain history and outstanding obligations
before the updated participant creates new intents.
Rollback may use only a binary compatible with the current persisted schema,
or a separately qualified reverse migration. Retain the latest journal and
chain evidence; restoring an older backup cannot erase newer obligations.

Coordinator deployments use the same versioned, resumable state discipline in
serverless storage and fence obsolete checkpoint and signing generations.
Migrations preserve its own chain journal, hotkey access, pending commitments,
and finalized cursor alongside cohort state. Operational
updates must preserve frozen evaluation meaning and unfinished work. A missing
migration or invalid signature pauses only its dependent actions with the
evidence retained. A chain runtime upgrade pauses nothing; it can only defer an
operation it leaves without a type-correct encoding. Release archives, migration
compatibility, and recovery after many offline versions are tested requirements.

Native validators use local journals; the coordinator retains its chain journal
alongside its application records in serverless storage. Monitor journal
capacity, retain completed evidence beyond bounded diagnostic history,
and qualify migration before exhaustion. Never delete operation identities or
remove journaling to bypass a refusal. A published correction keeps the
original result, records its authority and new revision, and accounts for
already-submitted chain effects. Changes to evaluation meaning belong in a
future comparison; an invalid current score requires an explicit correction
policy before rewards can proceed.
